\subsection{Tema 3. El grupo fundamental}
\indicetema{e3}

\begin{ejercicio}{T3.1 (la concatenación respeta $\simeq_{\{0,1\}}$)}
Sean $\alpha, \alpha' \in C_{x_0,x_1}(X)$ y $\beta, \beta' \in C_{x_1,x_2}(X)$ con $\alpha \simeq_{\{0,1\}} \alpha'$ y $\beta \simeq_{\{0,1\}} \beta'$. Prueba que $\alpha * \beta \simeq_{\{0,1\}} \alpha' * \beta'$. Observa que, de hecho, $\alpha * \beta \simeq_{\{0,\frac12,1\}} \alpha' * \beta'$.

\enapuntes{\cref{prop:ej_poincare_arcos}~a) (ejercicio de clase [C5.3]).} \quad \enexamen{final 29/01/2026, ej.\ 4a, versión con lazos (1 pt)}

\emph{Lo de $\{0,\frac12,1\}$} \fuentepropia: la homotopía de esa demostración cumple además $H(\frac12,t) = H^{(1)}(1,t) = x_1$ para todo $t$, así que también fija $s = \frac12$.
\end{ejercicio}

\begin{ejercicio}{T3.2 (inverso de una concatenación)}
Prueba que, si $\alpha$ y $\beta$ son arcos con $\alpha(1) = \beta(0)$, entonces $(\alpha * \beta)^{-1} = \beta^{-1} * \alpha^{-1}$.

\enapuntes{\cref{prop:ej_inverso_concat} (ejercicio de clase [C5.7]).}
\end{ejercicio}

\begin{ejercicio}{T3.3 (Poincaré para arcos)}
Utilizando el mismo argumento que en la prueba del teorema de Poincaré, prueba:
\begin{enumerate}
    \item[a)] Si $\alpha \in C_{x_0,x_1}(X)$, entonces $\alpha * c_{x_1} \simeq_{\{0,1\}} c_{x_0} * \alpha \simeq_{\{0,1\}} \alpha$.
    \item[b)] Si $\alpha \in C_{x_0,x_1}(X)$, entonces $\alpha * \alpha^{-1} \simeq_{\{0,1\}} c_{x_0}$ y $\alpha^{-1} * \alpha \simeq_{\{0,1\}} c_{x_1}$.
    \item[c)] Si $\alpha \in C_{x_0,x_1}$, $\beta \in C_{x_1,x_2}$ y $\gamma \in C_{x_2,x_3}$, entonces $\alpha * (\beta * \gamma) \simeq_{\{0,1\}} (\alpha * \beta) * \gamma$.
\end{enumerate}

\enapuntes{\cref{prop:ej_poincare_arcos} (ejercicio de clase [C5.3]). Correspondencia de apartados: a) del boletín es c) de los apuntes, b) es d) y c) es b).}
\end{ejercicio}

\begin{ejercicio}{T3.4 (arcos homotópicos inducen el mismo cambio de base)}
Sean $w, w' \in C_{x_0,x_1}(X)$. Prueba que, si $w \simeq_{\{0,1\}} w'$, entonces $P_w = P_{w'}$.

\enapuntes{\cref{prop:ej_pw_homotopicos} (ejercicio de clase [C5.12]).}
\end{ejercicio}

\begin{ejercicio}{T3.5 ($\pi_1$ abeliano y cambios de base)}
Sea $X$ arco-conexo. Prueba que $\pi_1(X,x)$ es abeliano si y solo si, para todo par de puntos $x_0, x_1 \in X$ y todo par de arcos $w, w' \in C_{x_0,x_1}(X)$, se cumple $P_w = P_{w'}$.

\sol \fuentepropia

\emph{Cálculo previo.} Sea $\delta = w^{-1} * w' \in C_{x_1}(X)$, un lazo en $x_1$. Para $[\alpha] \in \pi_1(X,x_0)$, metiendo $w * w^{-1} \simeq c_{x_0}$ (\textbf{T3.3} b) y usando \textbf{T3.1} y la asociatividad (T3.3 c):
\[ P_{w'}([\alpha]) = [w'^{-1} * w * w^{-1} * \alpha * w * w^{-1} * w'] = [\delta]^{-1} \cdot P_w([\alpha]) \cdot [\delta], \]
donde se ha usado $[w'^{-1} * w] = [(w^{-1} * w')^{-1}] = [\delta]^{-1}$ (\textbf{T3.2}).

$\Rightarrow$: como $X$ es arco-conexo, todos los $\pi_1(X,x)$ son isomorfos [C5.10], luego $\pi_1(X,x_1)$ es abeliano y $[\delta]^{-1} \cdot P_w([\alpha]) \cdot [\delta] = P_w([\alpha])$. Así $P_{w'} = P_w$.

$\Leftarrow$: sea $x_0 = x_1 = x$, $w = c_x$ y $w' = \gamma$ un lazo cualquiera en $x$. Entonces $P_{c_x} = 1$ (por T3.3 a) y $P_\gamma([\alpha]) = [\gamma]^{-1}[\alpha][\gamma]$. La hipótesis $P_\gamma = P_{c_x}$ da $[\gamma]^{-1}[\alpha][\gamma] = [\alpha]$, es decir, $[\alpha][\gamma] = [\gamma][\alpha]$ para todo $[\alpha], [\gamma]$: $\pi_1(X,x)$ es abeliano.
\end{ejercicio}

\begin{ejercicio}{T3.6 (simplemente conexo y arcos con los mismos extremos)}
Sea $X$ arco-conexo. Prueba que son equivalentes:
\begin{enumerate}
    \item[a)] $\pi_1(X)$ es el grupo trivial.
    \item[b)] Para todo par de puntos $x_0, x_1$ y todo par de arcos $\alpha, \beta$ de $x_0$ a $x_1$, se cumple $\alpha \simeq_{\{0,1\}} \beta$.
\end{enumerate}

\sol \fuentepropia

\emph{b)$\Rightarrow$a).} Dado un lazo $\alpha \in C_{x_0}(X)$, aplicamos b) con $x_1 = x_0$ y $\beta = c_{x_0}$: $[\alpha] = [c_{x_0}]$.

\emph{a)$\Rightarrow$b).} $\alpha * \beta^{-1}$ es un lazo en $x_0$; por a), $\alpha * \beta^{-1} \simeq_{\{0,1\}} c_{x_0}$. Usando \textbf{T3.1} y \textbf{T3.3}:
\[ \alpha \simeq \alpha * c_{x_1} \simeq \alpha * (\beta^{-1} * \beta) \simeq (\alpha * \beta^{-1}) * \beta \simeq c_{x_0} * \beta \simeq \beta, \]
todas relativas a $\{0,1\}$.
\end{ejercicio}

\begin{ejercicio}{T3.7 (grupo fundamental de un producto)}
Sean $X$ e $Y$ espacios topológicos y $(x_0,y_0) \in X \times Y$. Prueba que $\pi_1(X \times Y,(x_0,y_0)) \cong \pi_1(X,x_0) \times \pi_1(Y,y_0)$.

\enapuntes{\cref{prop:ej_pi1_producto} (ejercicio de clase [C9.8]).}

\emph{Consecuencias:} $\pi_1(T^2) = \pi_1(S^1 \times S^1) \cong \mathbb{Z}^2$ y $\pi_1(S^1 \times [0,1]) \cong \mathbb{Z} \times 1 \cong \mathbb{Z}$.
\end{ejercicio}

\begin{ejercicio}{T3.8 (producto por un contráctil)}
Sean $X$ e $Y$ espacios topológicos con $Y$ contráctil. Prueba que $X \times Y \simeq X$.

\sol \fuentepropia

Sea $F \colon 1_Y \simeq c_{y_0}$. Tomamos $f \colon X \times Y \to X$, $f(x,y) = x$, y $g \colon X \to X \times Y$, $g(x) = (x,y_0)$. Entonces $f \circ g = 1_X$ y $g \circ f(x,y) = (x,y_0)$. La aplicación
\[ H((x,y),t) = (x,\, F(y,t)) \]
es continua, con $H_0 = 1_{X \times Y}$ y $H_1 = g \circ f$. Luego $g \circ f \simeq 1$ y $X \times Y \simeq X$.
\end{ejercicio}

\begin{ejercicio}{T3.9 (quitar un punto en una equivalencia homotópica)}
Sea $f \colon X \to Y$ una equivalencia homotópica con $f^{-1}(f(p)) = \{p\}$. ¿Es cierto que la restricción $f| \colon X \setminus \{p\} \to Y \setminus \{f(p)\}$ es una equivalencia homotópica?

\sol \fuentepropia

\textbf{No.} Sea $Y = C = [-1,1] \times \{0\} \cup \{0\} \times [-1,1]$ la cruz y $X = \mathbb{R}$. Sea
\[ f \colon \mathbb{R} \to C, \qquad f(x) = \left(\frac{x}{1+|x|},\, 0\right). \]
$X$ es convexo y $C$ es estrellado respecto del origen, luego ambos son contráctiles. Así, \emph{cualquier} $f$ continua es una equivalencia homotópica: con $g \colon C \to \mathbb{R}$ cualquiera, $g \circ f \simeq 1_X$ y $f \circ g \simeq 1_C$ por \textbf{T2.6} c). Además $f^{-1}(f(0)) = f^{-1}(0,0) = \{0\}$.

Pero $f| \colon \mathbb{R} \setminus \{0\} \to C \setminus \{(0,0)\}$ no es una equivalencia homotópica. Supongamos que $g$ fuera una inversa homotópica y $H \colon f \circ g \simeq 1$ en $C \setminus \{(0,0)\}$. Sea $q = (0,\frac12)$, en el brazo superior. El arco $t \mapsto H(q,t)$ une $f(g(q))$, que está en el eje horizontal (la imagen de $f$), con $q$, que está en el brazo superior. Pero el brazo superior $\{0\} \times \left]0,1\right]$ es una arco-componente de $C \setminus \{(0,0)\}$ que no corta al eje horizontal. Contradicción.

(Con el Tema 5 es inmediato: $\#\pi_0(\mathbb{R} \setminus \{0\}) = 2 \neq 4 = \#\pi_0(C \setminus \{(0,0)\})$. Un ejemplo aún más simple: $f \colon \{p\} \to [0,1]$, $f(p) = 0$, es una equivalencia homotópica, pero $\emptyset \to \left]0,1\right]$ no puede serlo, porque no hay aplicaciones $\left]0,1\right] \to \emptyset$.)
\end{ejercicio}

\begin{ejercicio}{T3.10 (homotópica a constante $\Rightarrow$ $f_*$ nulo)}
Sea $f \colon X \to Y$ continua, $y_0 \in Y$ y $\gamma_{y_0} \colon X \to Y$ la constante $x \mapsto y_0$. Prueba que, si $f \simeq \gamma_{y_0}$, entonces $f_* \colon \pi_1(X,x) \to \pi_1(Y,f(x))$ es el homomorfismo nulo, para todo $x \in X$.

\sol \fuenteexamen{parcial 31/10/2025, ej.\ 3} \quad \enexamen{parcial 31/10/2025, ej.\ 3, literal (1,5 pts)}

Sea $x_0 \in X$ y $F \colon f \simeq \gamma_{y_0}$. Sea $\sigma \colon I \to Y$, $\sigma(t) = F(x_0,t)$, que va de $f(x_0)$ a $y_0$. Por el teorema de las aplicaciones homotópicas [C6.7], el diagrama
\[ \begin{array}{ccc} \pi_1(X,x_0) & \xrightarrow{\ f_{*,x_0}\ } & \pi_1(Y,f(x_0)) \\ & \searrow{\scriptstyle (\gamma_{y_0})_{*,x_0}} & \downarrow{\scriptstyle P_\sigma} \\ & & \pi_1(Y,y_0) \end{array} \]
conmuta, es decir, $P_\sigma \circ f_{*,x_0} = (\gamma_{y_0})_{*,x_0}$.

Para $\alpha \in C_{x_0}(X)$, $(\gamma_{y_0})_{*,x_0}([\alpha]) = [\gamma_{y_0} \circ \alpha] = [c_{y_0}]$, luego $(\gamma_{y_0})_{*,x_0}$ es el morfismo nulo. Por tanto, para todo $[\alpha]$,
\[ P_\sigma\big(f_{*,x_0}([\alpha])\big) = [c_{y_0}], \]
es decir, $f_{*,x_0}([\alpha]) \in \ker(P_\sigma) = \{[c_{f(x_0)}]\}$, porque $P_\sigma$ es isomorfismo. Así $f_{*,x_0}$ es el homomorfismo nulo. (Alternativamente: $f_{*,x_0} = P_\sigma^{-1} \circ (\gamma_{y_0})_{*,x_0}$.)

\textbf{Resultados utilizados} (el examen pide enunciarlos):
\begin{enumerate}
    \item[(a)] \emph{[C6.7]} Sean $f,g \colon X \to Y$ continuas, $F \colon f \simeq g$ y $x_0 \in X$. Sea $\sigma(t) = F(x_0,t)$ y $P_\sigma \colon \pi_1(Y,f(x_0)) \to \pi_1(Y,g(x_0))$ el isomorfismo inducido por $\sigma$. Entonces $g_{*,x_0} = P_\sigma \circ f_{*,x_0}$.
    \item[(b)] \emph{[C5.9]} Si $w$ es un arco de $x_0$ a $x_1$, la aplicación $P_w \colon \pi_1(X,x_0) \to \pi_1(X,x_1)$, $P_w([\alpha]) = [w^{-1} * \alpha * w]$, es un isomorfismo de grupos.
\end{enumerate}

\emph{Error típico:} no basta decir ``$f \circ \alpha \simeq \gamma_{y_0} \circ \alpha$''. Esa homotopía es libre (el punto base recorre $\sigma$) y $\gamma_{y_0} \circ \alpha$ es un lazo en otro punto base.
\end{ejercicio}
